% GENERATED FILE. Edit canonical lessons, then run npm run generate:books.
% canonical-source-hash: fnv1a64:4aa222df3b88af21
% canonical-lessons: LA-C540-pupa, LA-C540-trochus, LA-C540-turbo, LA-C540-harpastum, LA-C540-natatio

\chapter{Toys and play}
\label{ch:la-toys-and-play}
\hlchaptermodality{540}

\begin{chapteropening}
\textbf{By the end of this chapter:} \emph{I can talk about children's toys and games.}

Complete the last lesson of chapter 540: I can talk about children's toys and games.
\end{chapteropening}

\section[pūpa, pūpae]{pūpa, pūpae — a doll; a little girl}

\label{lesson:LA-C540-pupa}



\begin{quote}
Before the new one: say the Latin for play, then the Latin for to wrestle.
\end{quote}

\subsection*{You'll want to know: pūpa, pūpae}
\textbf{pūpa, pūpae} — \textquotedblleft{}a doll; a little girl\textquotedblright{}.

\begin{quote}
Fīliola cum pūpā suā dormit.
\end{quote}

\begin{quote}
\emph{The little daughter sleeps with her doll.}
\end{quote}

Feminine, first declension. Before her wedding a Roman girl gave her dolls to the household gods.

\begin{cousinweb}
\textbf{pūpa} is the feminine of \emph{pūpus} \textquotedblleft{}boy, child\textquotedblright{} (Lewis \& Short). English \emph{puppet} comes from it, and \emph{pupil} from the diminutive \emph{pūpillus}.
\end{cousinweb}

\begin{grammarlens}[title={What you've built}]
One more word for telling, explaining and describing at B1.
\end{grammarlens}

\subsection*{Your turn}
\begin{itemize}
\raggedright
  \item \emph{Say it:} \emph{pūpa, pūpae}
  \item \emph{Say it:} \emph{pūpa, pūpae}, once more
  \item \emph{Say it:} a sentence of your own with \emph{pūpa, pūpae}
  \item \emph{Recall:} say the Latin for play, then the Latin for to wrestle, then say \emph{pūpa, pūpae} again
\end{itemize}

\begin{tcolorbox}[breakable,colback=teal!4,colframe=teal!35!black,title={Before you move on}]
What does pūpa, pūpae mean? (\textquotedblleft{}A doll; a little girl\textquotedblright{}.) Say it once more.
\end{tcolorbox}
\section[trochus, trochī]{trochus, trochī — a hoop}

\label{lesson:LA-C540-trochus}



\begin{quote}
Before the new one: say the Latin for to wrestle, then the Latin for a doll.
\end{quote}

\subsection*{You'll want to know: trochus, trochī}
\textbf{trochus, trochī} — \textquotedblleft{}a hoop\textquotedblright{}.

\begin{quote}
Puer trochum per viam agit.
\end{quote}

\begin{quote}
\emph{The boy drives his hoop along the street.}
\end{quote}

Masculine, second declension. A metal hoop with little rings, driven with a stick.

\begin{cousinweb}
\textbf{trochus} was borrowed from Greek \emph{trochos} \textquotedblleft{}wheel, hoop\textquotedblright{} (Lewis \& Short; EtymDB).
\end{cousinweb}

\begin{grammarlens}[title={What you've built}]
One more word for telling, explaining and describing at B1.
\end{grammarlens}

\subsection*{Your turn}
\begin{itemize}
\raggedright
  \item \emph{Say it:} \emph{trochus, trochī}
  \item \emph{Say it:} \emph{trochus, trochī}, once more
  \item \emph{Say it:} a sentence of your own with \emph{trochus, trochī}
  \item \emph{Recall:} say the Latin for to wrestle, then the Latin for a doll, then say \emph{trochus, trochī} again
\end{itemize}

\begin{tcolorbox}[breakable,colback=teal!4,colframe=teal!35!black,title={Before you move on}]
What does trochus, trochī mean? (\textquotedblleft{}A hoop\textquotedblright{}.) Say it once more.
\end{tcolorbox}
\section[turbō, turbinis]{turbō, turbinis — a spinning top; a whirlwind}

\label{lesson:LA-C540-turbo}



\begin{quote}
Before the new one: say the Latin for a doll, then the Latin for a hoop.
\end{quote}

\subsection*{You'll want to know: turbō, turbinis}
\textbf{turbō, turbinis} — \textquotedblleft{}a spinning top; a whirlwind\textquotedblright{}.

\begin{quote}
Turbō in pavīmentō celeriter vertitur.
\end{quote}

\begin{quote}
\emph{The top spins fast on the floor.}
\end{quote}

Masculine, third declension. Anything that spins: a top, a whirlwind. English \emph{turbine}.

\begin{cousinweb}
\textbf{turbō} belongs with \emph{turba} \textquotedblleft{}commotion\textquotedblright{} and the verb \emph{turbō} \textquotedblleft{}to disturb\textquotedblright{} (Lewis \& Short; EtymDB): a thing in whirling motion. English \emph{turbine} comes from it.
\end{cousinweb}

\begin{grammarlens}[title={What you've built}]
One more word for telling, explaining and describing at B1.
\end{grammarlens}

\subsection*{Your turn}
\begin{itemize}
\raggedright
  \item \emph{Say it:} \emph{turbō, turbinis}
  \item \emph{Say it:} \emph{turbō, turbinis}, once more
  \item \emph{Say it:} a sentence of your own with \emph{turbō, turbinis}
  \item \emph{Recall:} say the Latin for a doll, then the Latin for a hoop, then say \emph{turbō, turbinis} again
\end{itemize}

\begin{tcolorbox}[breakable,colback=teal!4,colframe=teal!35!black,title={Before you move on}]
What does turbō, turbinis mean? (\textquotedblleft{}A spinning top; a whirlwind\textquotedblright{}.) Say it once more.
\end{tcolorbox}
\section[harpastum, harpastī]{harpastum, harpastī — a hand-ball game}

\label{lesson:LA-C540-harpastum}



\begin{quote}
Before the new one: say the Latin for a hoop, then the Latin for a spinning top.
\end{quote}

\subsection*{You'll want to know: harpastum, harpastī}
\textbf{harpastum, harpastī} — \textquotedblleft{}a hand-ball game\textquotedblright{}.

\begin{quote}
Iuvenēs in campō harpastō lūdunt.
\end{quote}

\begin{quote}
\emph{The young men play hand-ball on the field.}
\end{quote}

Neuter, second declension. A rough team game with a small ball, from Greek.

\begin{cousinweb}
\textbf{harpastum} was borrowed from Greek \emph{harpaston} (Lewis \& Short), a ball one snatches from another, from \emph{harpazo} \textquotedblleft{}to seize\textquotedblright{}.
\end{cousinweb}

\begin{grammarlens}[title={What you've built}]
One more word for telling, explaining and describing at B1.
\end{grammarlens}

\subsection*{Your turn}
\begin{itemize}
\raggedright
  \item \emph{Say it:} \emph{harpastum, harpastī}
  \item \emph{Say it:} \emph{harpastum, harpastī}, once more
  \item \emph{Say it:} a sentence of your own with \emph{harpastum, harpastī}
  \item \emph{Recall:} say the Latin for a hoop, then the Latin for a spinning top, then say \emph{harpastum, harpastī} again
\end{itemize}

\begin{tcolorbox}[breakable,colback=teal!4,colframe=teal!35!black,title={Before you move on}]
What does harpastum, harpastī mean? (\textquotedblleft{}A hand-ball game\textquotedblright{}.) Say it once more.
\end{tcolorbox}
\section[natātiō, natātiōnis]{natātiō, natātiōnis — swimming; a swimming place}

\label{lesson:LA-C540-natatio}



\begin{quote}
Before the new one: say the Latin for a spinning top, then the Latin for a hand-ball game.
\end{quote}

\subsection*{You'll want to know: natātiō, natātiōnis}
\textbf{natātiō, natātiōnis} — \textquotedblleft{}swimming; a swimming place\textquotedblright{}.

\begin{quote}
Natātiō in marī mihi placet.
\end{quote}

\begin{quote}
\emph{I like swimming in the sea.}
\end{quote}

Feminine, third declension, from \emph{natō} — to swim.

\begin{cousinweb}
\textbf{natātiō} is the action noun of \emph{natō} \textquotedblleft{}to swim\textquotedblright{} (EtymDB; Lewis \& Short). English \emph{natatorium}, a swimming pool, comes from the same verb.
\end{cousinweb}

\begin{grammarlens}[title={What you've built}]
One more word for telling, explaining and describing at B1.
\end{grammarlens}

\subsection*{Your turn}
\begin{itemize}
\raggedright
  \item \emph{Say it:} \emph{natātiō, natātiōnis}
  \item \emph{Say it:} \emph{natātiō, natātiōnis}, once more
  \item \emph{Say it:} a sentence of your own with \emph{natātiō, natātiōnis}
  \item \emph{Recall:} say the Latin for a spinning top, then the Latin for a hand-ball game, then say \emph{natātiō, natātiōnis} again
\end{itemize}

\begin{tcolorbox}[breakable,colback=teal!4,colframe=teal!35!black,title={Before you move on}]
What does natātiō, natātiōnis mean? (\textquotedblleft{}Swimming; a swimming place\textquotedblright{}.) Say it once more.
\end{tcolorbox}
